\subsection{The Third Sunday: \latin{Gaudete}}\label{sec:third-sunday}

In 2026 the Third Sunday falls on 13 December. The station is at St Peter's. The Mass opens \latin{Gaudete} (gloss: rejoice),
from the Epistle it will read, St Paul to the Philippians (4:4--7); the Gospel is John the
Baptist's testimony before the priests and Levites (John 1:19--28).

The day looks and sounds different. Under the Code of Rubrics of 1960, rose vestments may be
used on this Sunday, in its Office and Mass only (n.~131). The instruction of the Sacred
Congregation of Rites on sacred music of 1958 silences the organ and every other instrument in
the liturgy of Advent, from First Vespers of the First Sunday to None of the Vigil, except at
eucharistic Benediction (n.~81\,a). It makes exceptions. On this Sunday the organ or harmonium
may be played, and no other instrument (n.~83\,b); on the other days of the season the organ or
harmonium may be heard at Mass and Vespers only to sustain the singing (n.~83\,c); and on feasts
of precept and certain local feasts the organ and other instruments are allowed (n.~83\,a).

\subsubsection{The Lord is nigh}

The Epistle gives the command and its reason:

\begin{quote}
Rejoice in the Lord always: again, I say, rejoice.

Let your modesty be known to all men. The Lord is nigh.\par
{\smallskip\noindent\small ---Philippians 4:4--5, Douay-Rheims (Challoner).\par}
\end{quote}

\noindent Rejoicing here is a command, and in the next two verses Paul sets beside it the means of
obeying it: anxiety is to be turned into prayer made with thanksgiving, and the peace of God is
to keep the heart and mind in Christ Jesus.

The Collect asks for light. In Bute's English:

\begin{quote}
O Lord, we beseech Thee, mercifully incline Thine ears unto our prayers, and lighten the darkness
of our minds by the grace of Thy heavenly visitation; Who livest and reignest with God the Father,
in the unity of the Holy Ghost, one God, world without end.\par
{\smallskip\noindent\small ---Collect of the Third Sunday of Advent, \latin{Aurem tuam}, tr.
Bute, \textit{The Roman Breviary} (1908), Winter, p.~243.\par}
\end{quote}

The ``heavenly visitation'' the prayer asks for is a coming of God to the mind, which the
mind in its darkness cannot give itself. It is close to what Bernard, in his fifth sermon, calls
the hidden coming (Section~\ref{sec:three-comings}).

\subsubsection{The voice and the Word}

The Gospel is John's examination. The priests and Levites ask him whether he is the Christ, or
Elijah, or the prophet, and to each he says no. Then he says what he is:

\begin{quote}
He said: I am the voice of one crying in the wilderness, make straight the way of the Lord, as
said the prophet Isaias.\par
{\smallskip\noindent\small ---John 1:23, Douay-Rheims (Challoner).\par}
\end{quote}

\noindent He baptizes with water, and among them stands one they do not know, the latchet of whose
shoe he is not worthy to loose. The 1962 Office
reads at Matins the opening of St Gregory's homily on the passage, and Gregory begins with
John's humility:

\begin{quote}
Dearly beloved brethren, the first thing which striketh us in to-day's Gospel is the
lowly-mindedness of John. He was so great that it was thought he might be the Christ; yet he
soberly chose rather to seem only what he really was, than to let the belief of men invest him
with a dignity which did not belong to him; for ``he confessed, and denied not, but confessed, I
am not the Christ,'' at the same time he would not deny what he was in reality; and thus his very
truth-speaking made him a member of Him Whose title he would not by falsehood take. In that he
arrogated not to himself the name of Christ, he became a member of Christ. While he humbly strove
to confess his own weakness, he earned by his simplicity a part in the grandeur of his
Master.\par
{\smallskip\noindent\small ---Gregory the Great, \textit{Homiliae in Evangelia} 7.1, tr. Bute,
\textit{The Roman Breviary} (1908), Winter, p.~241.\par}
\end{quote}

Gregory then explains why John calls himself a voice. As the \textit{Catena Aurea} gives his
words:

\begin{quote}
Ye know that the only-begotten Son is called the Word of the Father. Now we know, in the case of
our own utterance, the voice first sounds, and then the word is heard. Thus John declares himself
to be the voice, i.e. because he precedes the Word, and, through his ministry, the Word of the
Father is heard by man.\par
{\smallskip\noindent\small ---Gregory the Great, \textit{Homiliae in Evangelia} 7.2, as the
\textit{Catena Aurea} gives it on John 1:23, Oxford tr. (1874), p.~49.\par}
\end{quote}

St Augustine, preaching on the same Gospel, notices what kind of herald John is:

\begin{quote}
Would it not have seemed to you that a herald would have cried, ``Go away, make room.'' Instead
of the herald's cry ``Go away,'' John says ``Come.'' The herald makes men stand back from the
judge; to the Judge John calls. Yes, indeed, John calls men to the lowly One, that they may not
experience what He will be as the exalted Judge.\par
{\smallskip\noindent\small ---Augustine, \textit{Tractates on John} IV.7, tr. John Gibb, NPNF,
1st ser., vol.~7 (New York, 1888).\par}
\end{quote}

In Augustine's last sentence the lowly One and the exalted Judge are the same Christ, and John
calls men to him now so that they need not meet him first as Judge. The Lord whom the Epistle
calls nigh comes first in lowliness.

\medskip
\noindent\textit{Counsel for the week.} Read Philippians 4:4--7 on the Sunday and do what it
says: name one anxiety and turn it into a request made with thanksgiving. Pray the Collect
\latin{Aurem tuam} each morning for light. Take John's examination as your own once this week:
what you are not, and what you are, before God, without inflation and without false modesty. The
Ember Days fall on Wednesday, Friday and Saturday (Section~\ref{sec:ember}), and the great O
antiphons begin on Thursday 17 December (Section~\ref{sec:last-days}); plan the week around them.
